\documentclass[14pt,a4paper,oneside]{extbook}
\usepackage[utf8]{vietnam}
\usepackage{amsmath,amssymb,mathrsfs,fancyhdr,enumerate,multirow,makecell}
\usepackage{tkz-euclide,tikz-3dplot,tikz,tkz-tab}
\usetikzlibrary{arrows.meta,calc,intersections,angles,quotes,backgrounds,shapes.geometric,patterns,through,shadows,positioning}
\usepackage{yhmath}
\usepackage[top=0.7cm, bottom=0.7cm, left=1.0cm, right=1.0cm, headheight=16pt, headsep=3mm, footskip=5mm]{geometry}
\usepackage[hidelinks,unicode]{hyperref}
\usepackage{qrcode}
\usepackage[many]{tcolorbox}
\usepackage{fontawesome}

\pagestyle{fancy}
\fancyhf{}
\fancyhead[L]{\footnotesize\bfseries\color{teal!85!black}\faBook~SỔ TAY KIẾN THỨC TOÁN 4}
\fancyhead[R]{\footnotesize\bfseries\color{blue!85!black}Lớp Toán Thầy Lĩnh -- loptoanthaylinh.vn}
\fancyfoot[C]{\footnotesize\bfseries\color{gray!80!black}-- Sổ tay kiến thức: Bài 08. Góc nhọn, góc tù, góc bẹt --}
\renewcommand{\headrulewidth}{0.4pt}

\begin{document}

% --- BANNER TIÊU ĐỀ BÀI HỌC ---
\noindent
\begin{tcolorbox}[
  enhanced,
  colback=teal!8!white,
  colframe=teal!80!black,
  arc=3mm, boxrule=1pt,
  left=3mm, right=3mm, top=2mm, bottom=2mm,
  halign=center
]
  {\Large\bfseries\color{teal!90!black}\faBookmark\ BÀI 08: GÓC NHỌN, GÓC TÙ, GÓC BẸT}\\[1.5pt]
  {\footnotesize\itshape\color{black!75} Nhận diện và so sánh 4 loại góc cơ bản dựa trên chuẩn góc vuông 90 độ}
\end{tcolorbox}

\vspace{2.5mm}

% --- KHỐI 1 & 2: LÝ THUYẾT & KĨ NĂNG ---
\noindent
\begin{minipage}[t]{0.49\textwidth}
  \begin{tcolorbox}[
    enhanced,
    title={\footnotesize\bfseries\color{white}\faShapes\ 1. Bốn loại góc cơ bản},
    colback=blue!2!white,
    colframe=blue!75!black,
    arc=2.5mm, boxrule=0.8pt, top=1.8mm, bottom=1.8mm, left=2mm, right=2mm
  ]
    \footnotesize
    $\bullet$ \textbf{Góc vuông:} Số đo bằng \textbf{$90^\circ$}.\\
    $\bullet$ \textbf{Góc nhọn:} Số đo \textbf{bé hơn $90^\circ$} ($0^\circ < \alpha < 90^\circ$).\\
    $\bullet$ \textbf{Góc tù:} Số đo \textbf{lớn hơn $90^\circ$ và bé hơn $180^\circ$} ($90^\circ < \alpha < 180^\circ$).\\
    $\bullet$ \textbf{Góc bẹt:} Số đo bằng \textbf{$180^\circ$} (bằng 2 góc vuông, hai cạnh tạo thành đường thẳng).
  \end{tcolorbox}
\end{minipage}%
\hfill
\begin{minipage}[t]{0.49\textwidth}
  \begin{tcolorbox}[
    enhanced,
    title={\footnotesize\bfseries\color{white}\faCheck\ 2. Dùng ê-ke để kiểm tra},
    colback=teal!2!white,
    colframe=teal!75!black,
    arc=2.5mm, boxrule=0.8pt, top=1.8mm, bottom=1.8mm, left=2mm, right=2mm
  ]
    \footnotesize
    $\bullet$ Đặt góc vuông của thước ê-ke trùng với đỉnh và 1 cạnh của góc cần kiểm tra:\\
    - Cạnh kia \textbf{nằm trong} góc ê-ke $\rightarrow$ \textbf{Góc nhọn}.\\
    - Cạnh kia \textbf{trùng} cạnh ê-ke $\rightarrow$ \textbf{Góc vuông}.\\
    - Cạnh kia \textbf{nằm ngoài} góc ê-ke $\rightarrow$ \textbf{Góc tù}.\\
    $\bullet$ Hai cạnh cùng nằm trên một đường thẳng $\rightarrow$ \textbf{Góc bẹt}.
  \end{tcolorbox}
\end{minipage}

\vspace{2.5mm}

% --- KHỐI 3: SƠ ĐỒ TƯ DUY ---
\noindent
\begin{tcolorbox}[
  enhanced,
  title={\footnotesize\bfseries\color{white}\faSitemap\ Sơ đồ tư duy: Phân loại 4 loại góc theo số đo chuẩn},
  colback=yellow!2!white,
  colframe=teal!85!black,
  arc=3mm, boxrule=0.9pt, top=1.8mm, bottom=1.8mm, left=1.5mm, right=1.5mm,
  halign=center
]
  \begin{tikzpicture}[
    scale=0.92,
    every node/.style={transform shape},
    boxN/.style={rectangle, rounded corners=4pt, fill=blue!10, draw=blue!75!black, text=blue!90!black, font=\scriptsize, inner sep=3pt, align=center, minimum width=2.6cm},
    boxV/.style={rectangle, rounded corners=4pt, fill=teal!15, draw=teal!80!black, text=teal!95!black, font=\scriptsize, inner sep=3pt, align=center, minimum width=2.6cm},
    boxT/.style={rectangle, rounded corners=4pt, fill=orange!12, draw=orange!75!black, text=orange!90!black, font=\scriptsize, inner sep=3pt, align=center, minimum width=2.6cm},
    boxB/.style={rectangle, rounded corners=4pt, fill=purple!12, draw=purple!75!black, text=purple!90!black, font=\scriptsize, inner sep=3pt, align=center, minimum width=2.6cm}
  ]
    % Hình vẽ góc nhọn
    \draw[thick, blue!80!black] (-5.1, 0.45) -- ++(35:0.75);
    \draw[thick, blue!80!black] (-5.1, 0.45) -- ++(0:0.85);
    \draw[blue!70!black] (-5.1+0.35, 0.45) arc (0:35:0.35);

    % Hình vẽ góc vuông
    \draw[thick, teal!85!black] (-1.7, 0.45) -- ++(90:0.75);
    \draw[thick, teal!85!black] (-1.7, 0.45) -- ++(0:0.85);
    \draw[teal!85!black] (-1.7, 0.45+0.25) -- ++(0.25,0) -- ++(0,-0.25);

    % Hình vẽ góc tù
    \draw[thick, orange!85!black] (1.7, 0.45) -- ++(130:0.75);
    \draw[thick, orange!85!black] (1.7, 0.45) -- ++(0:0.85);
    \draw[orange!80!black] (1.7+0.3, 0.45) arc (0:130:0.3);

    % Hình vẽ góc bẹt
    \draw[thick, purple!85!black] (5.1-0.85, 0.45) -- (5.1+0.85, 0.45);
    \fill[purple!85!black] (5.1, 0.45) circle (1.5pt);
    \draw[purple!80!black] (5.1+0.35, 0.45) arc (0:180:0.35);

    \node[boxN] (n1) at (-5.1, -0.35) {\textbf{GÓC NHỌN}\\[1pt]Số đo $< 90^\circ$};
    \node[boxV] (n2) at (-1.7, -0.35) {\textbf{GÓC VUÔNG}\\[1pt]Số đo $= 90^\circ$};
    \node[boxT] (n3) at (1.7, -0.35) {\textbf{GÓC TÙ}\\[1pt]$90^\circ < \text{Số đo} < 180^\circ$};
    \node[boxB] (n4) at (5.1, -0.35) {\textbf{GÓC BẸT}\\[1pt]Số đo $= 180^\circ$};

    \draw[teal!85!black, thick, -{Stealth[length=2mm]}] (n1.east) -- (n2.west);
    \draw[teal!85!black, thick, -{Stealth[length=2mm]}] (n2.east) -- (n3.west);
    \draw[teal!85!black, thick, -{Stealth[length=2mm]}] (n3.east) -- (n4.west);
  \end{tikzpicture}
\end{tcolorbox}

\vspace{2.5mm}

% --- KHỐI 4: MẸO & BẪY ---
\noindent
\begin{minipage}[t]{0.54\textwidth}
  \begin{tcolorbox}[
    enhanced,
    title={\footnotesize\bfseries\color{white}\faSortNumericAsc\ Mẹo so sánh độ lớn 4 loại góc},
    colback=green!3!white,
    colframe=green!60!black,
    arc=2.5mm, boxrule=0.8pt, top=1.8mm, bottom=1.8mm, left=2.2mm, right=2.2mm
  ]
    \footnotesize
    $\bullet$ \textbf{Thứ tự từ bé đến lớn:}\\
    $\textbf{Góc nhọn} < \textbf{Góc vuông} < \textbf{Góc tù} < \textbf{Góc bẹt}$\\
    $\bullet$ 1 góc bẹt $= 2$ góc vuông $= 180^\circ$.\\
    $\bullet$ 1 góc vuông $= 90^\circ$.\\
    $\bullet$ Góc nhọn tạo cảm giác "sắc nhọn", góc tù tạo cảm giác "mở rộng".
  \end{tcolorbox}
\end{minipage}%
\hfill
\begin{minipage}[t]{0.44\textwidth}
  \begin{tcolorbox}[
    enhanced,
    title={\footnotesize\bfseries\color{white}\faExclamationTriangle\ Mẹo tránh bẫy trắc nghiệm},
    colback=red!2!white,
    colframe=red!70!black,
    arc=2.5mm, boxrule=0.8pt, top=1.8mm, bottom=1.8mm, left=2.2mm, right=2.2mm
  ]
    \footnotesize
    $\bullet$ \textbf{Bẫy hình bị xoay nghiêng:}\\
    Khi hình vẽ không nằm ngang, mắt nhìn dễ nhầm góc nhọn thành góc vuông hoặc góc tù!\\
    $\rightarrow$ \textbf{Bắt buộc đặt thước ê-ke} để kiểm tra, không đoán mò bằng mắt!\\
    $\bullet$ Góc bẹt có 2 cạnh cùng nằm trên 1 đường thẳng.
  \end{tcolorbox}
\end{minipage}

\end{document}
